过去四年，医疗人工智能经历了三个层面的发展：训练范式上从任务专用监督学习走向通用基础模型，输出形式上从判别式判断走向生成式交互，模态范围上从单模态走向多模态融合。以 GPT-4、PaLM 2、Gemini 与 LLaMA 系列为代表的大语言与多模态模型，将医学问答、影像报告生成、临床决策支持与患者沟通纳入统一的生成式框架之下，并持续刷新 USMLE、MultiMedQA、MultiMedBench、MedHELM 等基准表现。本文以医疗人工智能的发展图景为主线，沿范式演进、医疗大语言模型、多模态医疗大模型与专科基础模型四条线索梳理当下医疗 AI 的能力基础（第2节至第5节），并在第6节讨论临床证据、评估方法与治理框架三方面目前仍未得到解决的若干具体问题。医疗 AI 的能力扩张由数据规模、训练目标与结构先验共同驱动，走出通用化与专科化彼此互补的技术路径。多种方法的层化叠加，构成了医疗人工智能当下能力的来源，也为后续讨论可信、可部署、可监管的实现路径奠定了基础。
